% Generated by scripts/gen_blueprint.py from blueprint/nodes/05-six-exponentials.toml. Do not edit.

\chapter{The six exponentials theorem}
\label{chap:05-six-exponentials}

Schneider's method in the form of Lang and Ramachandra, relativised to a fixed number field: an auxiliary function built by Siegel's lemma, a descent that extends its zeros over a lattice, and a zero estimate.

In this chapter $x=(x_1,\dots,x_d)$ and $y=(y_1,\dots,y_l)$ are families of complex numbers, each linearly independent over $\Q$. For $\lambda\in\N^{d}$ and $m\in\N^{l}$ write $\langle\lambda,x\rangle=\sum_i\lambda_ix_i$ and $\langle m,y\rangle=\sum_jm_jy_j$. For $L\in\N$ and integers $p_\lambda$ the auxiliary function is
\[ F_p(z)=\sum_{\lambda\in\{0,\dots,L-1\}^{d}}p_\lambda\,e^{\langle\lambda,x\rangle z}. \]
A number field $K$ is a subfield of $\C$ of finite degree over $\Q$.

% Prove2Me: SX.exists_aux_expSum -- https://prove2.me/theorems/f0a028dd-7409-4134-b87b-24b2921be5ae
\begin{lemma}[The auxiliary function]
\label{lem:exists_aux_expSum}
\lean{SX.exists_aux_expSum}
\leanok
Let $d+l<dl$, and let $K$ be a number field containing every $e^{x_iy_j}$. There is $c>0$ such that for every large $M$ there are an integer $L\ge1$ with $L^{d}\le cM^{l}$ and integers $p_\lambda$ with $|p_\lambda|\le e^{cLM}$, not all zero for $\lambda\in\{0,\dots,L-1\}^{d}$, such that $F_p(\langle m,y\rangle)=0$ for every $m\in\N^{l}$ with $m_j<M$ for all $j$.

\emph{Source:} \cite[Ch.~II]{Lang66}, \cite{Ram68}; see \cite[\S3.4]{Wal09}.
\end{lemma}

\begin{proof}
\leanok
\uses{lem:exists_denom_house_monomial_le,lem:siegel_entrywise}
At the lattice points the exponentials are products of powers of the $e^{x_iy_j}$, so the conditions form a linear system over $K$: $M^{l}$ conditions on $L^{d}$ unknowns, and $d+l<dl$ leaves room for Siegel's lemma (Lemma~\ref{lem:siegel_entrywise}), with the denominators and houses of Lemma~\ref{lem:exists_denom_house_monomial_le}.
\end{proof}

% Prove2Me: SX.descent_step -- https://prove2.me/theorems/d2ebac57-a209-470b-a366-8430f6fe9ff3
\begin{lemma}[The descent step]
\label{lem:descent_step}
\lean{SX.descent_step}
\leanok
Let $d+l<dl$, let $K$ be a number field containing every $e^{x_iy_j}$, and let $c>0$. There is $M_0$ such that for all $M\ge M_0$, all $L\ge1$ with $L^{d}\le cM^{l}$, all integers $p_\lambda$ with $|p_\lambda|\le e^{cLM}$ and all $N\ge M$: if $F_p(\langle m,y\rangle)=0$ for every $m\in\N^{l}$ with all $m_j<N$, then $F_p(\langle m,y\rangle)=0$ for every $m\in\N^{l}$ with all $m_j<N+1$.

\emph{Source:} \cite[Ch.~II]{Lang66}, \cite{Ram68}; see \cite[\S3.4]{Wal09}.
\end{lemma}

\begin{proof}
\leanok
\uses{lem:expPoly_grid_estimate,lem:exists_denom_house_monomial_le,lem:liouville_house}
At a new lattice point $F_p$ is small, by Lemma~\ref{lem:expPoly_grid_estimate} applied to the zeros already known. Its value there is an element of $K$ with controlled denominator and house (Lemma~\ref{lem:exists_denom_house_monomial_le}), so if it were non-zero Lemma~\ref{lem:liouville_house} would bound it from below. The two bounds are incompatible because $dl>d+l$.
\end{proof}

% Prove2Me: SX.eq_zero_of_expSum_vanishes -- https://prove2.me/theorems/fd939533-ac90-4615-9614-02acc8274f65
\begin{lemma}[Vanishing on the whole lattice]
\label{lem:eq_zero_of_expSum_vanishes}
\lean{SX.eq_zero_of_expSum_vanishes}
\leanok
Let $l\ge2$, $L\in\N$ and $p_\lambda$ integers. If $F_p(\langle m,y\rangle)=0$ for every $m\in\N^{l}$, then $p_\lambda=0$ for every $\lambda\in\{0,\dots,L-1\}^{d}$.

\emph{Source:} \cite[Ch.~II]{Lang66}, \cite{Ram68}; see \cite[\S3.4]{Wal09}.
\end{lemma}

\begin{proof}
\leanok
The exponents $\langle\lambda,x\rangle$ are pairwise distinct, and as $l\ge2$ and the $y_j$ are independent, $\lambda\mapsto(e^{\langle\lambda,x\rangle y_j})_j$ is injective; Dedekind's independence of characters concludes.
\end{proof}

% Prove2Me: SX.six_exponentials_of_numberField -- https://prove2.me/theorems/8e1e71c6-a4e1-49e6-83e3-41cc6024100d
\begin{theorem}[Six exponentials over a number field]
\label{thm:six_exponentials_of_numberField}
\lean{SX.six_exponentials_of_numberField}
\leanok
Let $d,l\in\N$ with $d+l<dl$, let $x_1,\dots,x_d$ and $y_1,\dots,y_l$ be complex numbers, each family linearly independent over $\Q$, and let $K$ be a number field. Then some $e^{x_iy_j}$ does not lie in $K$.

\emph{Source:} \cite[Ch.~II]{Lang66}, \cite{Ram68}; see \cite[Theorem 1.12]{Wal00}.
\end{theorem}

\begin{proof}
\leanok
\uses{lem:exists_aux_expSum,lem:descent_step,lem:eq_zero_of_expSum_vanishes}
Otherwise Lemma~\ref{lem:exists_aux_expSum} gives an auxiliary function, Lemma~\ref{lem:descent_step} makes it vanish on the whole lattice by induction, and Lemma~\ref{lem:eq_zero_of_expSum_vanishes} makes all its coefficients zero ($d+l<dl$ forces $l\ge2$), a contradiction.
\end{proof}

% Prove2Me: DiazModulus.six_exponentials -- https://prove2.me/theorems/7d362030-8eff-4aa9-8138-7b9b15cda0be
% Prove2Me: Schanuel.six_exponentials -- https://prove2.me/theorems/75ca528d-83e0-43c7-b183-42e42df38b0f
\begin{theorem}[The six exponentials theorem]
\label{thm:six_exponentials}
\lean{Diaz.six_exponentials}
\leanok
Let $x_1,x_2$ be complex numbers linearly independent over $\Q$, and let $y_1,y_2,y_3$ be complex numbers linearly independent over $\Q$. Then at least one of the six numbers $e^{x_iy_j}$ ($i=1,2$, $j=1,2,3$) is transcendental.

\emph{Source:} Lang \cite[Ch.~II]{Lang66}, Ramachandra \cite{Ram68}; see \cite[\S1.3, Theorem 1.12]{Wal00}.
\end{theorem}

\begin{proof}
\leanok
\uses{thm:six_exponentials_of_numberField}
Finitely many algebraic numbers lie in one number field; apply Theorem~\ref{thm:six_exponentials_of_numberField} with $(d,l)=(2,3)$.
\end{proof}
